\scanpage{62}
% manuscript line 1
\[
\boxed{(a_1\wedge\cdots\wedge a_n,b_1\wedge\cdots\wedge b_n)=\det\big((a_i,b_j)_{i,j\leq n}\big)} \tag*{(4)}
\]

% manuscript line 1a
% slanted, fragment on the folded corner
\margin{(a) \illegible{2} (\unsure{b|3})}

% manuscript line 2
On peut donc énoncer la première partie du

% manuscript line 3
\deleted{Prop\unsure{.}} \uline{Théorème} Soit \deleted{\unsure{un}} $E$ un espace de Hilbert. Il existe sur
$(\overset{n}{\bigwedge}E)\times(\overset{n}{\bigwedge}E)$ une forme hermitienne définie positive (et une seule)
satisfaisant à 4. Le complété est la \deleted{produit tensoriel} \inserted{puissance extérieure} hilbertienne
$n$-ième de $\overset{n}{\bigwedge}E$. \deleted{Si} Si $(e_i)_{i\in I}$ est une base \deleted{(\unsure{hilbert}} orthonormée
de $E$, alors les $e_J$ donnés par (5) forment une base
\deleted{\unsure{hilbert. Notons que} \illegible{2} $(e_J)$ \illegible{3}} hilbertienne
\deleted{\unsure{Nous allons se} \illegible{2}} dans $\overset{n}{\bigwedge}\vphantom{\bigwedge}^{(2)}E$. Enfin, on a
\deleted{\unsure{on a} $\unsure{\det((a_i,a_j))}$}

% manuscript line 11
\[
\boxed{\|a_1\wedge\cdots\wedge a_n\|=(\det((a_i,a_j)))^{1/2}\leq\|a_1\|\cdots\|a_n\|} \tag*{(5)}
\]

% manuscript line 12
\deleted{En effet}, l'égalité \inserted{dans (5)} est un cas particulier de (4) (pour $b_i=a_i$).
Considérons, dans \deleted{l'espace engendré par les} \inserted{un espace de dim. $n$ contenant les} $a_i$, \deleted{une}
base orthonormale $(e_i)$, \deleted{telle que l'\unsure{on ait} \illegible{2}}
\deleted{vérifie aussitôt que} \inserted{posons} $a_{\deleted{\unsure{\ldots}}\inserted{\textcircled{i}}}=\sum_j\deleted{a_{\unsure{ij}}}a_{ij}e_j$,
\deleted{\unsure{soit}} $A=(a_{ij})_{i,j\leq n}$, alors $(a_i,a_j)=\sum_k a_{ik}\overline{a_{jk}}$
\deleted{$\det((a_i,a_j))=\det$}
est l'élément d'indice $i,j$ de \inserted{la matrice} $AA^*$, donc
$\det((a_i,a_j))=\det AA^*=\det A\det A^*=\inserted{\det A\,\overline{\det A}=}|\det A|^2$, \deleted{d'où} il
reste à prouver \fbox{(6) $\lvert\det A\rvert\leq\|a_1\|\cdots\|a_n\|$} (inégalité
de Hadamard \unsure{!}), qui exprime que le volume du
parallélépipède \unsure{rectangle|construit} sur les $a_i$ est au plus
égal au produit des longueurs des arêtes).

% manuscript line 15a
% slanted
\margin{\unsure{Ajouter au théorème|Ajouter un théorème} pour \unsure{préambule} \illegible{2} que $X\to x\wedge X$ de $\overset{n}{\bigwedge}\vphantom{\bigwedge}^{(2)}E$ dans $\overset{n+1}{\bigwedge}\vphantom{\bigwedge}^{(2)}E$ est de norme $\leq\|x\|$.}

% manuscript line 24
Pour ceci, \inserted{il suffit de remarquer qu'}on peut supposer les $(e_i)$ choisis de façon
que $A$ soit une matrice triangulaire, donc
$\det A=a_{11}\,a_{22}\cdots a_{nn}$, ce qui achève la démonstration.

% manuscript line 27
\uline{Corollaire 1} Soient $a_i,b_i$ $(i=1,\ldots,n)$ des éléments d'un espace
de Hilbert $E$. Alors on a $|\det((a_i,b_j))|\leq\prod_i\|a_i\|\,\|b_i\|$

% manuscript line 29
% archive sheet number: 63
